% Optional examiner-style review annotations.
% Define \ThesisReviewMode before loading this file to show the annotations.
\newif\ifThesisReview
\ifdefined\ThesisReviewMode
  \ThesisReviewtrue
\else
  \ThesisReviewfalse
\fi

% Assessment colours: applied to the section that follows a heading.
\definecolor{reviewgreen}{RGB}{210,238,210}
\definecolor{reviewyellow}{RGB}{255,239,153}
\definecolor{reviewred}{RGB}{248,194,194}
\definecolor{reviewgrey}{RGB}{220,220,220}
% Status colours: applied to the text itself, not to an assessment of it.
\definecolor{reviewpurple}{RGB}{223,205,240}
\definecolor{reviewblue}{RGB}{170,214,247}
\definecolor{reviewbluemark}{RGB}{18,84,150}

\newcommand{\ReviewLabelgreen}{GREEN --- sufficient}
\newcommand{\ReviewLabelyellow}{YELLOW --- weak or unclear}
\newcommand{\ReviewLabelred}{RED --- machine-like prose pattern}
\newcommand{\ReviewLabelgrey}{GREY --- logical or evidential gap}
\newcommand{\ReviewLabelpurple}{PURPLE --- revised and confirmed}
\newcommand{\ReviewLabelblue}{BLUE --- addition}

\newcommand{\ReviewMark}[2]{%
  \ifThesisReview
    \par\smallskip
    \begingroup
      \setlength{\fboxsep}{5pt}%
      \noindent\colorbox{review#1}{%
        \parbox{\dimexpr\linewidth-2\fboxsep\relax}{%
          \small\textbf{\csname ReviewLabel#1\endcsname.} #2%
        }%
      }%
    \endgroup
    \par\smallskip
  \fi
}

\newcommand{\ReviewLegendItem}[3]{%
  \noindent\colorbox{review#1}{\strut\textbf{#2}}\quad #3\par\smallskip
}

% ------------------------------------------------------------------
% Inline span markers.
%
% \Revised{...} marks text that has been read, checked and settled. The wording
% inside it is final. It is not edited again, not reworded, and not extended:
% whatever the passage says is what it is to say.
%
% \Comment{...} is how a settled passage is questioned instead of changed. It
% prints as comment: "..." and is the only thing that may appear inside a
% revised passage. It is a remark about the text, not part of it, so it is
% suppressed entirely outside review mode.
%
% \Added{...} marks genuinely new thesis text, in a passage that has not been
% revised yet. It prints its argument unchanged outside review mode. It is not
% used inside a revised passage; a comment is used there instead.
%
% Two restrictions follow from how the highlight is typeset. It is laid out
% character by character, so anything it must not take apart has to be handed
% to it as a single box:
%   1. Inline maths, citations, cross-references and control spaces inside a
%      marked span are written inside \mbox{}: \mbox{\(J^\top\)},
%      \mbox{\citep{Key}}, \mbox{\Cref{fig:x}}, \mbox{et al.\ }. A macro that
%      hides the \mbox does not help, since the span is scanned before the
%      macro is expanded.
%   2. A span may not cross a paragraph break. Mark each paragraph separately.
% ------------------------------------------------------------------
\ifThesisReview
  \usepackage{soul}
  % Long code identifiers are split into boxed segments joined by \allowbreak so
  % that a line can break inside them. The highlighter has to be told to step
  % over the command rather than try to typeset it.
  \soulregister{\allowbreak}{0}
  % Distinguishes an addition from a revised passage where the two tints meet
  % in one sentence, and survives greyscale printing. Redefine to \relax to
  % suppress it.
  \newcommand{\ReviewAddedTag}{\textsuperscript{\textcolor{reviewbluemark}{+}}}
  \newcommand{\Revised}[1]{{\sethlcolor{reviewpurple}\hl{#1}}}
  \newcommand{\Added}[1]{{\sethlcolor{reviewblue}\hl{#1}}\ReviewAddedTag}
  \newcommand{\Comment}[1]{%
    {\sethlcolor{reviewblue}\hl{\mbox{\textbf{comment:}} ``#1''}}%
  }
\else
  \newcommand{\ReviewAddedTag}{}
  \newcommand{\Revised}[1]{#1}
  \newcommand{\Added}[1]{#1}
  % A comment is a remark about the thesis, never part of it.
  \newcommand{\Comment}[1]{}
\fi
